\documentclass[10pt]{article}

\usepackage[utf8]{inputenc}

\usepackage[T1]{fontenc}

\usepackage{amsmath}

\usepackage{amsfonts}

\usepackage{amssymb}

\usepackage[version=4]{mhchem}

\usepackage{stmaryrd}

\usepackage{graphicx}

\usepackage[export]{adjustbox}

\graphicspath{ {./images/} }

\usepackage{bbold}

\usepackage{CJKutf8}



\begin{document}

можно проиллюстрировать тем, что т. н. великая теорема Ферма допускает следующую әквивалентную формулировку: минимизировать



\[

\left(x_{1}^{t}+x_{2}^{t}-x_{3}^{t}\right)^{2}

\]



при условиях



\[

x_{1} \geqslant 1, x_{2} \geqslant 1, x_{3} \geqslant 1, t \geqslant 3,

\]



$t, x_{1}, x_{2}, x_{3}$ - целые числа. Если какой-то метод Ц. п. в качестве ответа выдаст положительное значение для минимума целевой функции, это будет конструктивным доказательством теоремы Ферма, если же ответом будет нуль, то это опровергнет ее.



Центральный теоретич. вопрос в Ц. п.: можно ли исключить полный перебор при решении задач Ц. п.? Одна из математич. формулировок этой проблемы: совпадают ли классы $P$ и $N P$ ? Класс $P$ (класс $N P$ ) - это переборные задачи, разрешимые на детерминированной (недетерминированной) машине Тьюринга за полиномиальное время, т. е. за число вычислительных операций, зависящее как полином от т. н. «длины записи» задачи. Класс $N P$ вклютает все задачи Ц. п., к-рые имеют экспоненциальное (относительно «длины записи» задачи) количество допустимых репений. Проблема «P=NP?» пока (1984) остается открытой.



\includegraphics[max width=\textwidth, center]{2023_07_09_605286dbbd033bba59c5g-1(1)}

направления в линейном программировании, М., 1966; [2]



\includegraphics[max width=\textwidth, center]{2023_07_09_605286dbbd033bba59c5g-1(2)}

Вычислительные машины и труднорешаемые задачи, пер. с англ., М., 1982.



ЦЕЛЫЙ ИДЕАЛ - идеал поля $Q$ относительно кольца $A$ (здесь $Q-$ поле частных кольца $A$ ), целиком лежащий в $A$. При этом Ц. и. является идеалом в $A$ и обратно, всякий идеал кольца $A$ - Ц. и. его поля частных $Q$.



ЦЕЛЫХ ТОЧЕК РАСПРЕДЕЛЕНИЕ - Нек-рые асимптотические формулы аналитич. теории чисел для арифметич. функций, к-рые могут быть сформулированы как задачи о числе целых точек в нек-рых многообразиях, в первую очередь, в гомотетически расширяющихся областях в пространстве $\mathbb{R} n$.



Классическими (исходными) здесь являются нруса проблема (Гаусса) и делителей проблема (Дирихле), а также их многочисленные обобщения.



Jum.: [1] F r i c k e r F., Einführung in die Gitterpunktlehre, Basel - Boston - Stuttgart, 1981; [2] X y a J o-r e H, \begin{CJK}{UTF8}{mj}年\end{CJK} сумм и его применения в теории чи-



\begin{center}

\includegraphics[max width=\textwidth]{2023_07_09_605286dbbd033bba59c5g-1}

\end{center}



ЦЕНТР - тиІ расположения траекторий автономной системы обыкновенных дифференциальных уравнений 2-го порядка



\[

\dot{x}=f(x), x=\left(x_{1}, x_{2}\right), f: G \subset \mathbb{R}^{2} \rightarrow \mathbb{R}^{2},

\]



$f \in C(G), G$ - область единственности, в окрестности особой точки $x_{0}$. Этот тип характеризуется следующцм образом. Существует окрестностк $U$ точки $x_{0}$ такая, что все траектории системы, начинаппцеся в $U \backslash\left\{x_{0}\right\}$, являются замкнутыми кривыми, окружаюццими $x_{0}$. Ц. наз. при этом и сама точка $x_{0}$. На рис. точка



\includegraphics[max width=\textwidth, center]{2023_07_09_605286dbbd033bba59c5g-1(3)}

$O$-центр. Движение по траекториям с возрастанием $t$ может происходить по часової стрелке или против нее (как указано на рис. стрелками). Ц. устойчив по Јяпунову (не асимптотически). Его индекс ІІанкаре равен 1.



Точка $x_{0}$ является для системы (\textit{) Ц., нашр., когда $f(x)=A\left(x-x_{0}\right)$, где $A-$ постоннная матрица с чисто мнимыми собственными значениямг. В отличие от простых точек гокоя других тиџов, встретаюџихся для линейных систем 2-го порядка (седло, узел, бокус), точка $x_{0}$ типа Ц. при возмущении линейной системы добавками к правой части, вообще говоря, не остается Ц., как бы ни были высоки порядок малости возмущений относительно $\left\|x-x_{0}\right\|$ и порядок их гладкости. Она может перейти при этом в фокус (устойчивый или неустойqивый) или в центрофокус (см. Центра и фокуса проблема). Для нелинейной системы (}) класса $C^{1}\left(f \in C^{1}(G)\right)$ точка покоя $x_{0}$ может быть Ц. и в том случае, когда матрица $A=f^{\prime}\left(x_{0}\right)$ имеет два нулевых собственных значения.



Лum.: [1] А м е лькин В. В., Лука ш е в ич Н. А., С а д о в с к и й А. П., Нелинейные колебания в системах второго порядка, Минск, 1982. См. также лит. при статье Особая $\begin{array}{ll}\text { точка дифференциального уравнения. } & \text { А. Ф. Андреeв. }\end{array}$



ЦЕНТР г у у п п - множество $Z$ всех ц е н тр а л ь в ы х (называемых иногда также и н в а p иа н т ны и ) әле ме н то в этой группы, т. е. әлементов, перестановочных со всеми элементами группы. Ц. группы $G$ является нормальным делителем в $G$ и даже характеристич. подгруппой в $G$. Более того, нормальным делителем в $G$ будет всякая подгруппа Ц. Абелевы групшы и только они совпадают со своим Ц. Группы, Ц. к-рых состоит лишь из единицы, наз. г р у п п а м и бе з це н т р а. Факторгрупша $G / Z$ группы $G$ по своему Ц. не обязана быть группой без Ц. Лит.: [1] К у р о ш А. Г., Теория групп, 3 изд., М., 1967.



IЕНТР к о л ь ц а - совокупность $Z$ всех әлементов кольца, перестановочных с любым элементом, т. е.



\[

Z=\{z \mid a z=z a \text { для всех } a\} .

\]



Ц. кольца оказывается подкольцом, содержащим вместе с каждым обратимым элементом әлемент, обратный к нему. Ц. кольца, являющегося алгеброй с единицей над полем, содержит основное поле (ср. Центральная алгебра). Л. А. Скорняков.



ЦЕНТР топологической динамической системь $\left\{S_{t}\right\}$ (потока или каскада с фазовым пространством $X$ ) $\frac{1}{X}$ наибольшее замкнутое инвариантное множество $A \subset X$, все точки к-рого являются неблуждающими точками для ограничения исходной системы на $A$. Ц. заведомо непуст, если лространство $X$ компактно (более общо, если имеется полутраектория с компактным замыканием). Дж. Гиркгоф, к-рый ввел понятие II., пользовался другим, но экві валентным, определением с помощыо нек-рого трансфинитного процесса (см. [1] - [4]). Число шагов последнего наз. г л у б и но й Ц. (фактически имеется несколько «глубив», так как этот процесс допускает нек-рые модификации). Глубина Ц. невелика для потоков на компактных многообразиях размерности $\leqslant 2$ (см. [5], [6]) и для каскадов, получающихся итерированием гомеоморфизма окружности или непрерывного отображения (даже необратимого) отрезка (см. [7]), но уже для потоков в $\mathbb{R}^{3}$ и на нек-рых открытых поверхностях может быть сколь угодно больгіим счетным трансфинитом (см. [8]-[10]). В полном метрич. пространстве Ц. совпадает с замыканием множества точек, обладающих свойством устойчивости по Пуассону.



Если $X$ - компакт, а $U$-. окрестность Ц., то траектория любой точіки $x \in X$ «троводит бо́льпую часть времени в $U »:$ доля на отрезке $[0, T]$ тех $t$, для к-рых $S_{t} x \in U$, стремится к 1 при $T \rightarrow \infty$. Впрочем, наименьшее замкнутое инвариантное множество, обладающее тем же свойством ( п р и тяже ни я, см. $[3], \mid 4])$, является, вообще говоря, только частью Ц.; в метризуемом случае оно совіадает с Замыканием объединения всех әргодицеских множеств.



Jum.: [1] B i rk h of f G. D., "Nachr. der Gesellschaft der Wiss. Göttingen. Niath.- phys.' Kl.», 1926, H. 1, S. $81-92$;



А 26 Математическая әнц., т. 5





\end{document}